% !TEX program = tectonic
% =============================================================================
%  AULA 9 — Redes Veiculares e Mobilidade Inteligente
%  Curso: Redes sem Fio  |  Profa. Anelise Munaretto
%  Conteúdo: C-V2X, 802.11p/ITS-G5, platooning, V2X safety.
%  Compilar: tectonic aula9.tex
% =============================================================================
\documentclass[aspectratio=169,11pt]{beamer}
% =============================================================================
%  PREÁMBULO COMÚN para as aulas de Redes sem Fio (Beamer + TikZ)
%  Incluido com % =============================================================================
%  PREÁMBULO COMÚN para as aulas de Redes sem Fio (Beamer + TikZ)
%  Incluido com % =============================================================================
%  PREÁMBULO COMÚN para as aulas de Redes sem Fio (Beamer + TikZ)
%  Incluido com \input{preamble.tex} en cada aula.
%  Paleta de cores e estilos são definidos aqui uma única vez.
% =============================================================================

\usepackage[T1]{fontenc}
\usepackage{amsmath,amssymb,mathtools}
\usepackage{graphicx}
\usepackage{tikz}
\usetikzlibrary{positioning,arrows.meta,calc,backgrounds,decorations.pathmorphing,fit,shadows,shapes.symbols,shapes.geometric}
\usepackage{pgfplots}
\pgfplotsset{compat=1.16}
\usepackage{tabularx,booktabs,multirow,array}
\usepackage[normalem]{ulem}
\usepackage{hyperref}

% ---------- Paleta de cores própria ----------
\definecolor{aneliseAzul}{RGB}{20,60,110}
\definecolor{aneliseVerde}{RGB}{30,140,70}
\definecolor{aneliseLaranja}{RGB}{215,120,20}
\definecolor{aneliseGris}{RGB}{90,90,90}

% ---------- Tema ----------
\usetheme{Madrid}
\usecolortheme{whale}
\setbeamercolor{structure}{fg=aneliseAzul}
\setbeamercolor{title}{fg=aneliseAzul}
\setbeamercolor{frametitle}{fg=aneliseAzul}
\setbeamercolor{block title}{fg=aneliseAzul,bg=aneliseGris!12!white}
\setbeamercolor{block body}{bg=aneliseGris!7!white}
\hypersetup{colorlinks=true,linkcolor=aneliseAzul,urlcolor=aneliseVerde}

% ---------- Comandos auxiliares ----------
\newcommand{\kurose}{Kurose \& Ross, \emph{Computer Networking: a Top-Down Approach}}
\newcommand{\forouzan}{Forouzan, \emph{Data Communications and Networking}}
\newcommand{\stallings}{Stallings, \emph{Wireless Communications \& Networks}}
\newcommand{\tanenbaum}{Tanenbaum, \emph{Computer Networks}}
\newcommand{\rfc}[1]{\href{https://www.rfc-editor.org/rfc/rfc#1.txt}{RFC~#1}}
% -----------------------------------------------------------------------------
%  Estilos TikZ comuns para desenhar redes (posicionamento relativo)
% -----------------------------------------------------------------------------
\tikzset{
  host/.style={draw,rounded corners=2mm,fill=aneliseAzul!15,inner sep=1.5mm,font=\scriptsize},
  rt/.style={draw,rounded corners=1mm,fill=aneliseLaranja!25,inner sep=1mm,font=\scriptsize},
  nube/.style={draw,cloud,aspect=2.2,fill=aneliseGris!15,inner sep=1.5mm,font=\scriptsize},
  el/.style={draw,rounded corners=1.2mm,fill=aneliseGris!18,inner sep=0.8mm,font=\scriptsize},
  enl/.style={thick},
  enlA/.style={thick,-Stealth},
} en cada aula.
%  Paleta de cores e estilos são definidos aqui uma única vez.
% =============================================================================

\usepackage[T1]{fontenc}
\usepackage{amsmath,amssymb,mathtools}
\usepackage{graphicx}
\usepackage{tikz}
\usetikzlibrary{positioning,arrows.meta,calc,backgrounds,decorations.pathmorphing,fit,shadows,shapes.symbols,shapes.geometric}
\usepackage{pgfplots}
\pgfplotsset{compat=1.16}
\usepackage{tabularx,booktabs,multirow,array}
\usepackage[normalem]{ulem}
\usepackage{hyperref}

% ---------- Paleta de cores própria ----------
\definecolor{aneliseAzul}{RGB}{20,60,110}
\definecolor{aneliseVerde}{RGB}{30,140,70}
\definecolor{aneliseLaranja}{RGB}{215,120,20}
\definecolor{aneliseGris}{RGB}{90,90,90}

% ---------- Tema ----------
\usetheme{Madrid}
\usecolortheme{whale}
\setbeamercolor{structure}{fg=aneliseAzul}
\setbeamercolor{title}{fg=aneliseAzul}
\setbeamercolor{frametitle}{fg=aneliseAzul}
\setbeamercolor{block title}{fg=aneliseAzul,bg=aneliseGris!12!white}
\setbeamercolor{block body}{bg=aneliseGris!7!white}
\hypersetup{colorlinks=true,linkcolor=aneliseAzul,urlcolor=aneliseVerde}

% ---------- Comandos auxiliares ----------
\newcommand{\kurose}{Kurose \& Ross, \emph{Computer Networking: a Top-Down Approach}}
\newcommand{\forouzan}{Forouzan, \emph{Data Communications and Networking}}
\newcommand{\stallings}{Stallings, \emph{Wireless Communications \& Networks}}
\newcommand{\tanenbaum}{Tanenbaum, \emph{Computer Networks}}
\newcommand{\rfc}[1]{\href{https://www.rfc-editor.org/rfc/rfc#1.txt}{RFC~#1}}
% -----------------------------------------------------------------------------
%  Estilos TikZ comuns para desenhar redes (posicionamento relativo)
% -----------------------------------------------------------------------------
\tikzset{
  host/.style={draw,rounded corners=2mm,fill=aneliseAzul!15,inner sep=1.5mm,font=\scriptsize},
  rt/.style={draw,rounded corners=1mm,fill=aneliseLaranja!25,inner sep=1mm,font=\scriptsize},
  nube/.style={draw,cloud,aspect=2.2,fill=aneliseGris!15,inner sep=1.5mm,font=\scriptsize},
  el/.style={draw,rounded corners=1.2mm,fill=aneliseGris!18,inner sep=0.8mm,font=\scriptsize},
  enl/.style={thick},
  enlA/.style={thick,-Stealth},
} en cada aula.
%  Paleta de cores e estilos são definidos aqui uma única vez.
% =============================================================================

\usepackage[T1]{fontenc}
\usepackage{amsmath,amssymb,mathtools}
\usepackage{graphicx}
\usepackage{tikz}
\usetikzlibrary{positioning,arrows.meta,calc,backgrounds,decorations.pathmorphing,fit,shadows,shapes.symbols,shapes.geometric}
\usepackage{pgfplots}
\pgfplotsset{compat=1.16}
\usepackage{tabularx,booktabs,multirow,array}
\usepackage[normalem]{ulem}
\usepackage{hyperref}

% ---------- Paleta de cores própria ----------
\definecolor{aneliseAzul}{RGB}{20,60,110}
\definecolor{aneliseVerde}{RGB}{30,140,70}
\definecolor{aneliseLaranja}{RGB}{215,120,20}
\definecolor{aneliseGris}{RGB}{90,90,90}

% ---------- Tema ----------
\usetheme{Madrid}
\usecolortheme{whale}
\setbeamercolor{structure}{fg=aneliseAzul}
\setbeamercolor{title}{fg=aneliseAzul}
\setbeamercolor{frametitle}{fg=aneliseAzul}
\setbeamercolor{block title}{fg=aneliseAzul,bg=aneliseGris!12!white}
\setbeamercolor{block body}{bg=aneliseGris!7!white}
\hypersetup{colorlinks=true,linkcolor=aneliseAzul,urlcolor=aneliseVerde}

% ---------- Comandos auxiliares ----------
\newcommand{\kurose}{Kurose \& Ross, \emph{Computer Networking: a Top-Down Approach}}
\newcommand{\forouzan}{Forouzan, \emph{Data Communications and Networking}}
\newcommand{\stallings}{Stallings, \emph{Wireless Communications \& Networks}}
\newcommand{\tanenbaum}{Tanenbaum, \emph{Computer Networks}}
\newcommand{\rfc}[1]{\href{https://www.rfc-editor.org/rfc/rfc#1.txt}{RFC~#1}}
% -----------------------------------------------------------------------------
%  Estilos TikZ comuns para desenhar redes (posicionamento relativo)
% -----------------------------------------------------------------------------
\tikzset{
  host/.style={draw,rounded corners=2mm,fill=aneliseAzul!15,inner sep=1.5mm,font=\scriptsize},
  rt/.style={draw,rounded corners=1mm,fill=aneliseLaranja!25,inner sep=1mm,font=\scriptsize},
  nube/.style={draw,cloud,aspect=2.2,fill=aneliseGris!15,inner sep=1.5mm,font=\scriptsize},
  el/.style={draw,rounded corners=1.2mm,fill=aneliseGris!18,inner sep=0.8mm,font=\scriptsize},
  enl/.style={thick},
  enlA/.style={thick,-Stealth},
}

\title[Redes sem Fio]{Redes sem Fio\\ Redes Veiculares e Mobilidade Inteligente}
\subtitle{Aula 9}
\author[Anelise Munaretto]{Profa. Anelise Munaretto}
\institute[Curso]{Redes sem Fio --- Ciência da Computação}
\date{2026}

\begin{document}

\begin{frame}[plain]
  \titlepage
\end{frame}

\section{Introdução}

\begin{frame}{Redes Veiculares: conceitos básicos}
  \begin{itemize}
    \item \textbf{VANET} (Vehicular Ad-hoc Network): rede formada por veículos e infraestrutura viária
    \item Comunicação nas modalidades:
      \begin{itemize}
        \item V2V (veículo-veículo)
        \item V2I (veículo-infraestrutura)
        \item V2P (veículo-pedestre)
        \item V2N (veículo-rede)
      \end{itemize}
    \item Aplicações:
      \begin{itemize}
        \item segurança (collision avoidance, emergency brake)
        \item eficiência (platooning, gestão de tráfego)
        \item conforto (infotainment, atualizações OTA)
      \end{itemize}
  \end{itemize}
\end{frame}

\section{PADRÕES}

\begin{frame}{Padrões: IEEE 802.11p vs.\ C-V2X}
  \begin{columns}
    \begin{column}{0.5\textwidth}
      \textbf{IEEE 802.11p / ITS-G5}
      \begin{itemize}
        \item derivado do Wi-Fi (OFDM, 10~MHz)
        \item banda 5,9~GHz (ITS)
        \item alcance típico ~300~m
        \item latência ~50~ms
        \item comunicação V2V/V2I direta
        \item sem infraestrutura celular
      \end{itemize}
    \end{column}
    \begin{column}{0.5\textwidth}
      \textbf{C-V2X (Cellular V2X)}
      \begin{itemize}
        \item baseado em LTE/5G NR (3GPP Rel.\ 14/16)
        \item modo \emph{PC5} (sidelink direto) e \emph{Uu} (via estação base)
        \item alcance $>$ 500~m (PC5)
        \item latência $<$ 20~ms (5G NR-V2X)
        \item suporte a \emph{sensing} cooperativo
      \end{itemize}
    \end{column}
  \end{columns}
  \vspace{0.2cm}
  \begin{exampleblock}{Tendência 2026}
    C-V2X (NR-V2X) domina novos projetos; IEEE 802.11p em operação em frotas existentes. China/Europa/América do Norte adotam C-V2X.
  \end{exampleblock}
\end{frame}

\begin{frame}{NR-V2X (3GPP Rel.\ 16/17)}
  \begin{itemize}
    \item Evolução do LTE-V2X usando 5G NR
    \item Características-chave:
      \begin{itemize}
        \item \textbf{sidelink} avançado: suporte a \emph{unicast}, \emph{groupcast}, \emph{broadcast}
        \item \textbf{modo 1}: escalonamento gerenciado pelo gNB
        \item \textbf{modo 2}: escalonamento autônomo (sensing de recursos)
      \end{itemize}
    \item Requisitos: latência 3--10~ms, confiabilidade 99,99\%
    \item Casos de uso: \emph{platooning}, \emph{advanced driving}, \emph{remote driving}
  \end{itemize}
  \begin{block}{C-V2X Day-1 vs.\ Day-2}
    \begin{itemize}
      \item \textbf{Day-1}: alertas básicos (freio, veículo de emergência, obras)
      \item \textbf{Day-2}: percepção cooperativa, intenção de manobra, \emph{trajectory sharing}
    \end{itemize}
  \end{block}
\end{frame}

\section{Platooning}

\begin{frame}{Platooning: comboio de veículos}
  \begin{columns}
    \begin{column}{0.55\textwidth}
      \begin{itemize}
        \item Veículos seguem em formação fechada (distância 5--15~m)
        \item \textbf{Líder} define velocidade; demais repetem automaticamente
        \item Benefícios:
          \begin{itemize}
            \item redução de arrasto aerodinâmico (até 20\% economia)
            \item menor congestionamento
            \item segurança: reação mais rápida que humana
          \end{itemize}
        \item Desafios:
          \begin{itemize}
            \item latência $<$ 10~ms
            \item falha de comunicação
            \item regulamentação
          \end{itemize}
      \end{itemize}
    \end{column}
    \begin{column}{0.45\textwidth}
      \begin{center}
        \begin{tikzpicture}[font=\scriptsize,
            v/.style={draw,rounded corners=1mm,fill=aneliseAzul!15,inner sep=1mm,text width=1.2cm,align=center},
          ]
          \node[v](l){Líder};
          \node[v,right=1.6cm of l](a){V1};
          \node[v,right=1.6cm of a](b){V2};
          \node[v,right=1.6cm of b](c){V3};
          \draw[thick,aneliseLaranja] (l.east) to[bend left=10] (a.west);
          \draw[thick,aneliseLaranja] (a.east) to[bend left=10] (b.west);
          \draw[thick,aneliseLaranja] (b.east) to[bend left=10] (c.west);
          \node[below=1mm of a,font=\tiny]{C-V2X sidelink};
        \end{tikzpicture}
      \end{center}
    \end{column}
  \end{columns}
\end{frame}

\section{Simulação}

\begin{frame}{Simulação de redes veiculares}
  \begin{itemize}
    \item \textbf{SUMO} (Simulation of Urban MObility): mobilidade veicular
    \item \textbf{ns-3 / OMNeT++}: rede de comunicação
    \item Framework acoplado: \textbf{Veins} (OMNeT++ + SUMO)
    \item C-V2X: \textbf{Simu5G}, \textbf{ns-3-mmwave}, \textbf{Sionna}
  \end{itemize}
  \begin{exampleblock}{Exercício proposto}
    Simular um cenário urbano com 50 veículos, um semáforo inteligente V2I e medir latência/pacotes perdidos. Comparar 802.11p vs.\ C-V2X.
  \end{exampleblock}
\end{frame}

\section{Bibliografia}

\begin{frame}{Bibliografia}
  \begin{itemize}
    \item 3GPP TR 37.885, \emph{Study on evaluation methodology of V2X}
    \item 3GPP TS 23.287, \emph{Architecture enhancements for V2X services}
    \item IEEE Std 802.11p-2010
    \item ETSI EN 302 571, \emph{ITS-G5}
    \item Sommer, Dressler, \emph{Vehicular Networking}, Cambridge Univ. Press, 2018
  \end{itemize}
\end{frame}

\end{document}